% Laget automatisk av elevark/lag_elevark.py fra opplegg/heis-mobil.html. Endringer her blir overskrevet.
\documentclass[11pt]{article}
\usepackage{elevark}
\begin{document}
\elevarkhode{Fysikk 1 · Krefter}{Heisen med mobilen}{https://fysikklab.no/opplegg/heis-mobil.html}{Opplegget på nett}

Når heisen setter i gang, kjenner du at du blir tyngre eller lettere. Med akselerometeret i mobilen kan dere måle akselerasjonen gjennom hele turen. Fra den kan dere regne ut farten, hvor høyt heisen gikk, og hva en badevekt ville vist.

\section*{Hensikt}

\begin{itemize}
\item Måle akselerasjonen i en heis og dele turen inn i faser med tilnærmet konstant akselerasjon.
\item Regne ut fart og høyde fra akselerasjonen, og sammenligne høyden med trykkmåleren i mobilen.
\item Bruke Newtons andre lov til å finne normalkraften på en person i heisen.
\end{itemize}

\section*{Utstyr}

\begin{itemize}
\item Mobil med appen phyphox (gratis) og eksperimentet for heis («Elevator» på engelsk)(Det er mulig at du ikke har sensorer i telefonen til dette. Da vil dette eksperimentet ikke være mulig og det vil stå inne i phyphox). Det viser akselerasjonen og høyden. Høyden regnes ut fra lufttrykket, så den virker bare på mobiler med trykkmåler (barometer). Uten trykkmåler bruker dere «Akselerasjon (uten $g$)».
\item En heis som går minst 3–4 etasjer, og gjerne en badevekt.
\end{itemize}

\section*{Framgangsmåte}

\begin{enumerate}
\item Legg mobilen flatt på gulvet i heisen med skjermen opp. Da måler z-aksen akselerasjonen opp og ned.
\item Start målingen, vent et par sekunder, og trykk på knappen for etasjen dere skal til. Hold mobilen i ro til heisen har stoppet helt, og vent et par sekunder til før dere stopper målingen.
\item Gjør gjerne en tur ned også, og en tur der én i gruppa står på badevekta og filmer skjermen på vekta.
\item Se på grafen over akselerasjonen. Del turen inn i faser: akselerasjon, jevn fart og oppbremsing. Les av gjennomsnittsakselerasjonen og varigheten til hver fase, og skriv dem inn i tabellen. Akselerasjon oppover er positiv.
\item Les av høydeforskjellen fra phyphox, eller tell etasjer og mål høyden på én etasje.
\end{enumerate}

\section*{Oppgaver}

\begin{oppgaveliste}
\oppgave \fase{gjett}{Gjett}Heisen går oppover. I hvilke deler av turen viser en badevekt mer enn vanlig, og når viser den mindre? Hva tror du vekta viser midt i turen, når farten er størst?
\svar
\oppgave \fase{test}{Mål}Gjør målingene og fyll inn tabellen. Ser fartsgrafen rimelig ut? Hva bør farten være når heisen har stoppet?
\svar
\oppgave \fase{test}{Regn}Hvor høyt gikk heisen ifølge akselerasjonsmålingene? Sammenlign med høyden fra trykkmåleren eller med antall etasjer.
\svar
\oppgave \fase{test}{Regn}Tegn kreftene på en person på 60~kg i hver fase. Regn ut normalkraften fra gulvet, og hva badevekta viser.
\svar
\oppgave \fase{forklar}{Forklar}Gjør målingen på vei ned også. Når viser badevekta mer enn vanlig nå? Forklar hvorfor oppstart oppover og oppbremsing nedover kjennes likt.
\svar
\oppgave \fase{forklar}{Vurder}Farten når heisen stopper blir sjelden nøyaktig null i utregningen. Hvorfor ikke? Hvilke andre feilkilder er det?
\svar
\oppgave \fase{utfordring}{Utfordring}Tenk deg at kabelen ryker og heisen faller fritt. Hva viser badevekta da? Hva viser «Akselerasjon (uten $g$)» og «Akselerasjon med $g$» i phyphox? Prøv å forklare hvorfor de to viser forskjellige tall.
\svar
\end{oppgaveliste}

\end{document}
